% !TEX root = ../main.tex
% (this is to make TexShop know where your parent file is)

% The text in the these appendix files is presented single spaced, if you want to do doublespacing,
% You need to change it to \textbackslash{doublespacing} % ie \doublespacing
% in the main template file main.tex

\setlength{\parskip}{0.2\baselineskip} % use this if yuo want to make a bit of space between paragraphs.
% !TEX root = ../main.tex
% (this is to make TexShop know where your parent file is)

%This is a template for tutor feedback on module. I tended to put it as an appendix, but you can copy paste it to any file.

\DefTblrTemplate{caption-tag}{default}{Feedback\hspace{0.25em}}
\DefTblrTemplate{caption-sep}{default}{:\enskip}
\DefTblrTemplate{caption-text}{default}{\InsertTblrText{caption}}


\phantomsection
\subsubsection*{René Boscio}
\addcontentsline{toc}{subsection}{René Boscio}

\begin{tblr}
       {
       colspec = {XXXX},
       width = \linewidth,
        vlines = {black},
        hlines = {black},
        cell{1,2}{odd} = {NiceGray},
        cell{1,2}{even} = {white},
        cell{2,4}{odd} = {NiceGray},
        cell{3}{1} = {r=1,c=4}{NiceGray},
        cell{4}{1} = {r=1,c=4}{c,white},
        cell{5}{1} = {r=1,c=4}{NiceGray},
        cell{6}{1} = {r=1,c=4}{white}
       }

    Date: & 17-04-2025 & Tutor: & René Boscio \\
    Module Code: & TSE701  & Via:     & 121 \\
  % \SetCell[c=4]{c} Feedback/notes from tutor:  \\
\end{tblr}


\begin{longtblr}
        {
        %hspan =  minimal,
        baseline = f,
        width = \linewidth,
        colspec = {X [1,l]},
        }
        Areas that were successful: Really good job recreating the sound of the original track. Listening to them one after the other, I could've easily believed I was listening to the same track again. From there on, you successfully expand the piece into a 2 minute main title-like cue, which great shape and electronic production, akin to modern contemporary electronic crime scores.

        Points for improvement: You strategically take out the low-end for the middle section, and we discussed how it goes on for a little too long without a low-frequency foundation. My recommendation was to gradually bring a new bass phrase (or pulse variation) after about 8 bars of being without it, to help amp up the stakes leading up to the next section.

        General advice: Overall I think this track was really well executed, and aside from the bass thing, my only other subjective advice was to ditch the bell chime sound at the end, and perhaps do a more intense swell that abruptly drops out which has become a bit of a stylistic trend within scores of this genre.

        Resources to investigate: I recommended checking out the End Credits (Extended Version) cue of The Watchful Eye score for production on the ending, and I would also recommend a few other scores such as the Theodor Shapiro's Severance Main Theme, and Ben Salisbury and Geoff Borrow's Archive 81 - Titles.

\end{longtblr}

\begin{tblr}
        {
        colspec = {X},
        width = \linewidth,
         vlines = {black},
         hlines = {black},
         cell{1}{odd} = {NiceGray},
        }


    I addressed this by: \\




   % \SetCell[c=4]{c} Feedback/notes from tutor:  \\
 \end{tblr}

 \begin{enumerate}
        \item Adding in a bass halfway through the middle section (as suggested) [01:21]
        \item Changing the ending to use the abrupt swell stop at the end (as in the reference tracks René mentioned) [02:00]

 \end{enumerate}

 \pagebreak

 \phantomsection
 \subsubsection*{Joel Grant}
 \addcontentsline{toc}{subsection}{Joel Grant}

 \begin{tblr}
        {
        colspec = {XXXX},
        width = \linewidth,
         vlines = {black},
         hlines = {black},
         cell{1,2}{odd} = {NiceGray},
         cell{1,2}{even} = {white},
         cell{2,4}{odd} = {NiceGray},
         cell{3}{1} = {r=1,c=4}{NiceGray},
         cell{4}{1} = {r=1,c=4}{c,white},
         cell{5}{1} = {r=1,c=4}{NiceGray},
         cell{6}{1} = {r=1,c=4}{white}
        }

     Date: & 28-04-2025 & Tutor: & Joel Grant \\
     Module Code: & TSE701  & Via:     & 121 \\
   % \SetCell[c=4]{c} Feedback/notes from tutor:  \\
 \end{tblr}


 \begin{longtblr}
         {
         %hspan =  minimal,
         baseline = f,
         width = \linewidth,
         colspec = {X [1,l]},
         }

         From memory (and short form notes), Joel told me that:
         \begin{enumerate}
            \item The opening bass was a little loud
            \item The bass in the continuation was very wide in the stereo image. It should be much narrower.
            \item I can push into the final limiter a little more to make the track louder overall.
            \item He liked what I'd done with the track, but that maybe I'd strayed too far from what Guy would have done. However, he emphasized that it was a good thing and advised me not to change it.
         \end{enumerate}

 \end{longtblr}

 \begin{tblr}
         {
         colspec = {X},
         width = \linewidth,
          vlines = {black},
          hlines = {black},
          cell{1}{odd} = {NiceGray},
         }


     I addressed this by: \\
    % \SetCell[c=4]{c} Feedback/notes from tutor:  \\
  \end{tblr}

  \begin{enumerate}
         \item Turning down the opening bass just a little [00:00--00:28]
         \item Fixing the stereo image. Turns out that the pulsing bass and the ping pong toms at the end were quite wide even below 200hz.
         \item Automating elements of the ``master'' chain, including how much gain is applied. I couldn't just set a static level, because it'd make the opening sound more different from Guy's version. Additionally, in automating, it had to be fairly subtle, such that any shift wouldn't be too obvious.\footnote{For the record, there are two places where the input gain is raised: [00:28] (instantaneous) and [00:52--01:00] (gradually over four bars)}
         \item As per Joel's suggestion, I did not change anything about the track.
  \end{enumerate}
